\documentclass[11pt]{article}
\usepackage[margin=1in]{geometry}
\usepackage[T1]{fontenc}
\usepackage{lmodern}
\usepackage{microtype}
\usepackage{xcolor}
\usepackage{hyperref}
\usepackage{booktabs}
\usepackage{enumitem}
\usepackage{listings}
\usepackage{tcolorbox}
\usepackage{parskip}
\usepackage{array}
\newcolumntype{L}[1]{>{\raggedright\arraybackslash}p{#1}}

% \path{} (from url, loaded by hyperref) prints paths and identifiers
% verbatim -- so `--` is not turned into a dash -- and may break at / _ - .
\makeatletter
\g@addto@macro\UrlBreaks{\do\_\do\-\do\.}
\makeatother
\urlstyle{tt}
\setlength{\emergencystretch}{3em}

\hypersetup{
  colorlinks=true,
  linkcolor=blue!60!black,
  urlcolor=blue!60!black
}

\definecolor{codebg}{RGB}{245,245,245}
\definecolor{codecomment}{RGB}{90,100,90}
\definecolor{codekeyword}{RGB}{0,70,140}
\definecolor{codestring}{RGB}{150,60,20}

\lstdefinestyle{pythonstyle}{
  language=Python,
  backgroundcolor=\color{codebg},
  basicstyle=\ttfamily\small,
  keywordstyle=\color{codekeyword}\bfseries,
  commentstyle=\color{codecomment}\itshape,
  stringstyle=\color{codestring},
  showstringspaces=false,
  breaklines=true,
  frame=single,
  rulecolor=\color{black!20},
  columns=fullflexible,
  keepspaces=true
}
\lstdefinestyle{shellstyle}{
  language=bash,
  backgroundcolor=\color{codebg},
  basicstyle=\ttfamily\small,
  showstringspaces=false,
  breaklines=true,
  frame=single,
  rulecolor=\color{black!20},
  columns=fullflexible,
  keepspaces=true
}
\lstset{style=pythonstyle}

\newtcolorbox{keyidea}{
  colback=blue!5, colframe=blue!50!black,
  title=Key Idea, fonttitle=\bfseries
}
\newtcolorbox{tipbox}{
  colback=green!5, colframe=green!45!black,
  title=Tip, fonttitle=\bfseries
}
\newtcolorbox{warningbox}{
  colback=orange!7, colframe=orange!60!black,
  title=Common Pitfall, fonttitle=\bfseries
}

\title{
  \textbf{Codebase Change Summary --- Week 5}\\[4pt]
  \large Mock vs.\ MagicMock; Control-Flow Graphs and Coverage\\[2pt]
  \normalsize SE 413/513 --- Software Testing
}
\author{Demo Testing Suite}
\date{}

\begin{document}
\maketitle

\begin{center}
\textit{What changed in \path{demo-testing-suite/} since Week 4, and why}
\end{center}

\vspace{0.3cm}

% ============================================================
\section{Purpose of This Document}
% ============================================================
The \path{demo-testing-suite/} codebase is a single, shared project that
grows one lecture at a time rather than being rebuilt from scratch each
week. This document summarizes exactly what was added or changed for
Week 5 and, more importantly, explains \emph{why} each change exists ---
which testing concept from the two Week 5 readings it demonstrates:
\begin{itemize}
\item \textit{Mock vs.\ MagicMock in Python} (closing out the test-doubles unit), and
\item \textit{Control-Flow Graphs and Statement, Branch, and Path Coverage}.
\end{itemize}
Read this alongside the readings and \path{demo-testing-suite/README.md}.

\begin{keyidea}
Nothing from Weeks 1--4 was removed or rewritten. Everything below is
\emph{additive}: the running examples you already saw are all still
there, unchanged, including their two intentional \texttt{xfail} results
and one intentional \texttt{FAIL}.
\end{keyidea}

% ============================================================
\section{What Changed At a Glance}
% ============================================================
\begin{center}
\small
\begin{tabular}{@{}L{5.9cm}L{1.5cm}L{7.0cm}@{}}
\toprule
\textbf{File} & \textbf{Status} & \textbf{Purpose} \\
\midrule
\multicolumn{3}{@{}l}{\textit{Reading 1: Mock vs.\ MagicMock}}\\
\path{app/accounts/user_lookup.py} & New & The reading's Examples
2--6, verbatim: functions that use their dependency through
\texttt{len()}, indexing, \texttt{in}, iteration, and \texttt{with}. \\
\path{tests/accounts/test_user_lookup.py} & New & Ten tests: why a
plain \texttt{Mock} fails, each MagicMock example, and the
single-use-iterator pitfall. \\
\midrule
\multicolumn{3}{@{}l}{\textit{Reading 2: Control-Flow Graphs and Coverage}}\\
\path{app/payments/transactions.py} & New & The reading's two CFG
examples: \path{classify_transaction()} (no loops, 4 paths) and
\texttt{contains\_negative()} (a loop with a back edge). \\
\path{tests/payments/test_transactions.py} & New & One test per path
(\(V(G)=4\)), plus the three basis paths (\(V(G)=3\)) and a ``many''
case for the loop. \\
\path{app/pricing/membership.py} & New &
\path{apply_membership_discount()}: the statement-vs.-branch
coverage example. \\
\path{tests/pricing/test_membership.py} & New & Both branches: the
test that gives 100\% statement coverage, plus the one branch coverage
additionally demands. \\
\path{demos/coverage_gaps/} & New & Deliberately under-tested code
whose \texttt{pytest-cov} reports reproduce the reading: the
\texttt{2->4} branch gap, and a Python ``goto fail'' analogue. \\
\midrule
\multicolumn{3}{@{}l}{\textit{Tooling and documentation}}\\
\texttt{requirements.txt} & Modified & Adds \texttt{pytest-cov}. \\
\path{docs/DEMO_GUIDE_week05_mocking_coverage.md} & New &
Section-by-section lecture companion for both readings. \\
\texttt{README.md} (suite and repository) & Modified & Updated layout,
coverage command, expected output, and a Weekly Updates index. \\
\bottomrule
\end{tabular}
\end{center}

\begin{tipbox}
Install the new dependency before running anything this week:
\begin{lstlisting}[style=shellstyle]
cd demo-testing-suite
pip install -r requirements.txt
\end{lstlisting}
\end{tipbox}

% ============================================================
\section{Part 1: Mock vs.\ MagicMock}
% ============================================================
\textbf{Locations:} \path{app/accounts/user_lookup.py},
\path{tests/accounts/test_user_lookup.py}

Every double you wrote in Week 4 was used only through \textbf{ordinary
method calls} (\path{repository.insert_user(...)},
\path{payment_service.process_payment(100)}), so \texttt{Mock} was
always enough. This week's app file is the first code in the suite
whose dependencies are used through \textbf{magic methods}:

\begin{center}
\footnotesize
\begin{tabular}{@{}lll@{}}
\toprule
\textbf{Function} & \textbf{Operation} & \textbf{Magic method(s)} \\
\midrule
\path{process_users(repository)} & \texttt{len(repository)}, \texttt{repository[0]} & \texttt{\_\_len\_\_}, \texttt{\_\_getitem\_\_} \\
\texttt{check\_permission(user, p)} & \texttt{p in user.permissions} & \texttt{\_\_contains\_\_} \\
\texttt{get\_names(users)} & \texttt{for user in users} & \texttt{\_\_iter\_\_} \\
\texttt{read\_file(file)} & \texttt{with file as f} & \texttt{\_\_enter\_\_}, \texttt{\_\_exit\_\_} \\
\texttt{get\_user\_name(connection, id)} & \texttt{with connection.cursor() as c} & \texttt{\_\_enter\_\_}, \texttt{\_\_exit\_\_} \\
\bottomrule
\end{tabular}
\end{center}

The test file walks the reading's examples in order. Three tests are
worth singling out:

\begin{itemize}
\item \path{test_plain_mock_raises_type_error_on_len} and
\path{test_plain_mock_works_once_the_magic_method_is_attached_by_hand}
show the reading's core difference directly: a plain \texttt{Mock}
fails at \texttt{len()}, until you attach \texttt{\_\_len\_\_} yourself,
which is what \texttt{MagicMock} does for you.
\item \path{test_check_permission_mixes_mock_and_magicmock} uses
a \texttt{Mock} for \texttt{user} and a \texttt{MagicMock} for
\texttt{user.permissions} in the same test: the choice is made
\emph{per object}, not per test.
\item \path{test_iter_return_value_is_silently_empty_the_second_time}
\textbf{passes}, and that is the pitfall from the reading's Example 4
warning box: configuring \texttt{iter([...])} means the second loop
quietly sees nothing.
\end{itemize}

\begin{lstlisting}[style=shellstyle]
pytest tests/accounts/test_user_lookup.py -v
\end{lstlisting}

\begin{tipbox}
Compare this file with Week 4's
\path{tests/accounts/test_registration_service.py}. Both mock
collaborators of account code, but only one needs \texttt{MagicMock}.
Before reading either test, look at the \emph{production} code and
predict which. That is the reading's closing Key Idea: \textbf{let the
code under test choose the mock.}
\end{tipbox}

% ============================================================
\section{Part 2: Control-Flow Graphs and Path Coverage}
% ============================================================
\textbf{Locations:} \path{app/payments/transactions.py},
\path{tests/payments/test_transactions.py}

\path{app/payments/transactions.py} holds the reading's two CFG
examples, verbatim, so you can draw each graph from real code and then
check your hand-computed numbers against a tool. The tests are
parametrized with ids that name the path each case exercises:

\begin{center}
\footnotesize
\begin{tabular}{@{}llcL{7.6cm}@{}}
\toprule
\textbf{Function} & \textbf{Decisions} & \(V(G)\) & \textbf{Test ids} \\
\midrule
\texttt{classify\_transaction} & 3 & 4 & \texttt{amount<0}, \texttt{intl-under-1000},\\
 & & & \texttt{domestic}, \texttt{intl-over-1000} \\
\texttt{contains\_negative} & 2 & 3 & \texttt{zero-iterations}, \path{one-iteration-early-return},\\
 & & & \path{one-iteration-loop-exits}, \path{many-iterations-negative-late} \\
\bottomrule
\end{tabular}
\end{center}

\begin{lstlisting}[style=shellstyle]
pytest tests/payments/test_transactions.py -v \
    --cov=app.payments.transactions --cov-branch --cov-report=term-missing
\end{lstlisting}

The report shows 100\% branch coverage. For
\path{classify_transaction()}, which has no loops, the four tests
also achieve 100\% \emph{path} coverage. For \texttt{contains\_negative()}
they cannot: the loop's back edge makes the number of paths grow with
the input length, so the tests follow basis-path testing instead.

% ============================================================
\section{Part 2: Statement vs.\ Branch Coverage}
% ============================================================
\textbf{Locations:} \path{app/pricing/membership.py},
\path{tests/pricing/test_membership.py},
\path{demos/coverage_gaps/myapp.py},
\path{demos/coverage_gaps/test_myapp.py}

\path{apply_membership_discount()} appears twice, on purpose:
\begin{itemize}
\item In the main suite (\path{app/pricing/membership.py}), with both
branches tested, so it reaches 100\% branch coverage.
\item In \path{demos/coverage_gaps/myapp.py}, as a four-line copy
with no docstring, so its line numbers match the reading. Its test
file covers \emph{only} the member branch.
\end{itemize}

Run the demo with and without branch measurement:
\begin{lstlisting}[style=shellstyle]
pytest demos/coverage_gaps/test_myapp.py --cov=myapp --cov-report=term-missing
pytest demos/coverage_gaps/test_myapp.py --cov=myapp --cov-branch --cov-report=term-missing
\end{lstlisting}
The first run reports \textbf{100\%}. The second reproduces the
reading's report exactly:
\begin{lstlisting}[style=shellstyle]
Name                           Stmts   Miss Branch BrPart  Cover   Missing
--------------------------------------------------------------------------
demos/coverage_gaps/myapp.py       4      0      2      1    83%   2->4
\end{lstlisting}

\begin{warningbox}
Statement coverage alone reported 100\% for a function whose
\texttt{False} branch was never executed. Always pass
\path{--cov-branch}.
\end{warningbox}

% ============================================================
\section{Part 2: A Python ``goto fail''}
% ============================================================
\textbf{Locations:} \path{demos/coverage_gaps/handshake.py},
\path{demos/coverage_gaps/test_handshake.py}

The reading's case study is Apple's CVE-2014-1266, where a duplicated
line made the signature check unreachable. \texttt{handshake.py} ships a
Python version of the same mistake: a copy-pasted \texttt{if} that lost
its \texttt{not}, so the params and signature checks can never run.

\begin{lstlisting}
def verify_handshake(hash_ok, params_ok, signature_ok):
    if not hash_ok:
        return False
    if hash_ok:  # KNOWN BUG -- copy-paste of the check above, lost its `not`
        return True
    if not params_ok:
        return False
    if not signature_ok:
        return False
    return True
\end{lstlisting}

Both tests in \texttt{test\_handshake.py} pass. Coverage tells a
different story:
\begin{lstlisting}[style=shellstyle]
pytest demos/coverage_gaps/test_handshake.py \
    --cov=handshake --cov-report=term-missing
\end{lstlisting}
reports only 50\% of statements executed, with lines \textbf{30--34} ---
every check after the bug --- missing.

\begin{keyidea}
Why a copy-pasted \texttt{if}, and not simply a stray
\texttt{return True}? Code after an \emph{unconditional}
\texttt{return} is statically dead. Python drops it when compiling, and
\texttt{coverage.py} then does not count it as a statement at all, so
the report would say 100\%. Coverage reports can only show you lines
that \emph{could} run; to show the skipped checks, the bypass has to be
conditional, as the real \texttt{goto fail} was.
\end{keyidea}

% ============================================================
\section{Part 2: What Coverage Cannot Tell You}
% ============================================================
No new code: this reuses Week 3's \path{determine_shipping_cost()}.
\begin{lstlisting}[style=shellstyle]
pytest tests/pricing/test_shipping.py \
    --cov=app.pricing.shipping --cov-branch --cov-report=term-missing
\end{lstlisting}
The report shows \textbf{100\% branch coverage}, and in the same output
\path{test_valid_weights[5.0-10.0]} is still \texttt{XFAIL}. Coverage
cannot see that 5.0 is the boundary that exposes the \texttt{< 5} bug;
only boundary-value analysis found it.

% ============================================================
\section{Coverage on the Whole Suite}\label{sec:whole-suite}
% ============================================================
Pointing the tool at the whole codebase is a good habit, and this
week it finds something:
\begin{lstlisting}[style=shellstyle]
pytest --cov=app --cov-branch --cov-report=term-missing
\end{lstlisting}
Every Week 5 module is at 100\%. Two Week 4 files are not, for opposite
reasons:
\begin{itemize}
\item \path{app/orders/order_service.py}: \textbf{80\%, line 37
missing}. A \emph{real} gap: no test drives \texttt{place\_order()} down
its \texttt{"Payment failed"} branch. It is left in place for you to
close (see Suggested Next Steps).
\item \path{app/notifications/password_reset.py}: \textbf{63\%}, with
\texttt{SmtpEmailClient} missing. A \emph{deliberate} gap: that class is
the real SMTP transport, and testing it would need a real mail server,
which is why the Week 4 seam exists. Chasing 100\% here would make the
suite worse, not better.
\end{itemize}

% ============================================================
\section{Updated: Expected Test Results}
% ============================================================
From the project root:
\begin{lstlisting}[style=shellstyle]
cd demo-testing-suite
pip install -r requirements.txt
pytest -v
\end{lstlisting}
you should now see \textbf{65 passed, 2 xfailed, 1 failed}: twenty
more passing tests than Week 4's baseline, and the same two
\texttt{xfail} results and one \texttt{FAIL} as before, unrelated to this
week's additions. \path{demos/coverage_gaps/} is not part of that
run (\texttt{pytest.ini} sets \texttt{testpaths = tests}); run it by path.

% ============================================================
\section{How This Connects to the Readings}
% ============================================================
\begin{keyidea}
Both readings ask you to look at the code under test before writing
the test. The mocking reading says to let the production code's
\emph{interface} choose the double: ordinary calls mean \texttt{Mock};
\texttt{len}, indexing, \texttt{in}, iteration, and \texttt{with} mean
\texttt{MagicMock}. The coverage reading says to let the production
code's \emph{structure}, its control-flow graph, show you which tests
are missing, and then to remember that structure is blind to values.
This week's files let you practice both against real, running code.
\end{keyidea}

% ============================================================
\section{Suggested Next Steps}
% ============================================================
\begin{enumerate}
\item Before opening \path{tests/accounts/test_user_lookup.py}, read
each function in \path{app/accounts/user_lookup.py} and decide,
per object, whether you would use \texttt{Mock} or \texttt{MagicMock}.
\item Draw the CFG of \texttt{contains\_negative()} and compute
\(V(G) = E - N + 2\). Check that it matches decisions \(+ 1 = 3\).
\item Close the real gap: add a test to
\path{tests/orders/test_order_service.py} whose payment stub returns
\texttt{False}, then re-run the whole-suite coverage command from
Section~\ref{sec:whole-suite} and confirm \texttt{order\_service.py} reaches 100\%.
\item In \path{demos/coverage_gaps/test_handshake.py}, add the test
\texttt{verify\_handshake(True, True, False) is False} and watch it
fail. Then fix \texttt{handshake.py}.
\item Try the readings' paper exercises (\texttt{products},
\texttt{load\_config}, \texttt{grade\_letter}) first by hand, then
write them as code and use \path{--cov-branch} to check your
hand-computed coverage.
\end{enumerate}

\end{document}
